\chapter{Прототип системы}\label{ch:6}

Реализован веб-прототип, принимающий аудиофайл (MP3, WAV, FLAC, OGG), выполняющий анализ через AST v2 + post-processing и отображающий результат в синхронизированном дашборде из шести панелей. Архитектура~--- трёхзвенная: Python-сервис инференса, TypeScript-backend, React-фронтенд.

\begin{figure}[H]
\centering
\includegraphics[width=0.6\textwidth]{fig_landing.png}
\caption{Главная страница прототипа: editorial hero, карточки <<В~системе>> / <<Архитектура>> / <<Лучшие метрики>> и загрузка трека.}
\label{fig:landing}
\end{figure}

\section{Архитектура веб-приложения}\label{sec:6-1}

Прототип состоит из трёх сервисов, соединённых REST-API:

\begin{lstlisting}
+----------------+    HTTP    +-----------------+   HTTP    +-----------------+
|  Frontend      | ---------> |  Backend        | --------> |  ML API         |
|  React + Vite  |            |  Elysia + Bun   |           |  FastAPI        |
|  port 5173     | <--------- |  port 3000      | <-------- |  port 8000      |
+----------------+            +-----------------+           +-----------------+
\end{lstlisting}

\textbf{ML API} (FastAPI, port~8000): AST~v2 в памяти (574\,МБ), endpoints \code{POST /analyze}, \code{POST /spectrogram}, \code{GET /health}. \textbf{Backend} (Elysia на Bun, port~3000): \code{POST/GET /api/tracks}, \code{POST /api/tracks/:id/analyze}, \code{GET /api/tracks/:id/audio|spectrogram}. Хранение~--- JSON в \code{web/backend/data/}, аудио в \code{uploads/}. \textbf{Frontend} (React~19 + Vite, port~5173)~--- SPA с маршрутами \code{/} и \code{/tracks/:id}. Разработка~--- Bun-скрипт \code{scripts/dev-all.ts}, запускающий все три сервиса параллельно.

\section{Синхронизированный дашборд}\label{sec:6-2}

\begin{figure}[H]
\centering
\includegraphics[width=0.6\textwidth]{fig_dashboard_blues.png}
\caption{Дашборд трека <<The Thrill Is Gone>> (B.\,B.~King \& Muddy Waters): плеер, A/V-trajectory, жанр, мел-спектрограмма, кривые arousal/valence.}
\label{fig:dashboard-blues}
\end{figure}

Главное решение~--- синхронизированный дашборд из шести панелей с общим playhead: курсор движется по всем панелям, hover подсвечивает момент, клик ставит pinned-закладку.

\textbf{Фреймворк синхронизации.} \code{DashboardContext} (\code{playheadTime}, \code{hoverTime}, \code{pinnedTime}, \code{isPlaying}, \code{seek}/\code{togglePlay}) с rAF 60\,fps; хук \code{useTimeScale} на \code{scaleLinear} из d3-scale + \code{ResizeObserver}; обёртка \code{ChartFrame}. Поверх каждой панели~--- \code{TimeCursor} с Playhead (красный solid+glow), Pinned (синий+label) и Hover (dashed).

\begin{lstlisting}[language=JavaScript]
const ts = useTimeScale(containerRef, paddingLeft, paddingRight);
const pixelX = ts.scale(timeInSeconds);     // time -> pixel
const time = ts.invert(pixelXFromMouse);    // pixel -> time
\end{lstlisting}

\subsection{Шесть аналитических панелей}

\textbf{Структура (\code{StructurePanel}).} Canvas-полоска с цветовым делением + список с временами и confidence; клик~--- seek и pin.

\textbf{Эмоциональный анализ (\code{EmotionPanel}).} Две полоски-таймлайна arousal/valence + canvas-кривая регрессионных значений.

\textbf{Эмоциональная траектория (\code{AVTrajectory}).} Квадрат 320$\times$320 (Russell circumplex) с 4~квадрантами, серой траекторией за весь трек и красным «кометом» из последних 20\,с. Координаты рассчитываются из классификационных вероятностей:
\begin{equation}
x_{\text{val}} = 0.5 + 0.5 \cdot (P(\text{Bright}) - P(\text{Dark})), \quad
y_{\text{aro}} = 0.5 + 0.5 \cdot (P(\text{High}) - P(\text{Low})).
\end{equation}
Регрессионные предсказания смещены к центру~--- все треки скапливались бы в «Грустное».

\textbf{Жанр (\code{GenrePanel}).} Top-3 с прогресс-барами; warning при top-1 $<$ 45\,\%; dominant genre timeline; stacked area 10 жанров. Адресует OOD (§7.3).

\textbf{Мел-спектрограмма (\code{SpectrogramPanel}).} Canvas 128$\times$T с viridis и gamma $\gamma=0{,}4$. Endpoint \code{/spectrogram} возвращает downsampled JSON (max~512 frames). Академически ценна~--- показывает, \textbf{что «видит» модель}.

\textbf{Сводка (\code{SummaryPanel}).} Pie распределения сегментов; доминирующий жанр; темп/тональность (librosa); средняя энергия и настроение; число структурных переходов.

\section{Ключевые UX-решения}\label{sec:6-3}

\begin{itemize}
  \item \textbf{Top-3 жанров + confidence warning}~--- академическая честность к OOD.
  \item \textbf{Russell circumplex по классификации, не регрессии}~--- устраняет bias.
  \item \textbf{Display-name эвристика} (\code{displayName.ts})~--- DP-сегментация: \code{0375\_dynamite.mp3} $\to$ \code{Dynamite}, распознаёт GTZAN-формат \code{genre.NNNNN}.
  \item \textbf{Loading skeletons и empty states}~--- shimmer-placeholders снижают perceived latency.
  \item \textbf{Поиск и сортировка}~--- toolbar с search + sort (Recent/Название/Длительность/Жанр)~--- критично при 20+ треках.
  \item \textbf{Приглушённая палитра} (Tableau~10 muted): синий \code{\#2563eb} для interactive, красный \code{\#dc2626} для playhead.
\end{itemize}

\section{Технологический стек}\label{sec:6-5}

\begin{table}[H]
\centering
\caption{Технологический стек прототипа}
\begin{tabular}{L{3cm}L{4cm}L{7cm}}
\toprule
Слой & Технология & Обоснование \\
\midrule
ML inference & Python 3.12 + PyTorch + FastAPI & Прямое использование checkpoint из обучения \\
Backend & Bun + Elysia + TypeScript & Быстрый старт ($\sim$100\,мс), нативный TypeScript \\
Frontend & React 19 + Vite + TypeScript & Стандарт индустрии, hot reload \\
Визуализация & d3-scale + d3-shape + d3-scale-chromatic & Лёгкие модули для scales и colormaps \\
Стили & Plain CSS с CSS variables & Полный контроль над стилем \\
State & React Context (DashboardContext) & Достаточно для прототипа \\
\bottomrule
\end{tabular}
\end{table}

Отвергнуты: Redux/Zustand, TailwindCSS, Recharts/Visx, Node.js (Bun в 4~раза быстрее).

\begin{figure}[H]
\centering
\includegraphics[width=0.6\textwidth]{fig_tracklist.png}
\caption{Галерея треков: cover-tiles с ID3/YouTube-обложками, жанровые чипы, поиск и сортировка.}
\label{fig:tracklist}
\end{figure}

\section{Визуальный редизайн и расширения интерфейса}\label{sec:6-6}

\textbf{Dark Glass Theme.} Система переработана в стилистике \emph{glassmorphism}: полупрозрачные поверхности с \code{backdrop-filter: blur(24px) saturate(130\%)} на \code{\#0a0c14} с многослойными \code{radial-gradient} тинтами цвета настроения (\code{--track-mood}). Типографика~--- моноширинные uppercase-метки и серифный Fraunces.

\begin{figure}[H]
\centering
\includegraphics[width=0.6\textwidth]{fig_vibemode.png}
\caption{Режим Vibe~Mode: полноэкранный WebGL2-шейдер с SDF raymarching, реагирующий на FFT, баса/громкости envelope, onset и жанровые пресеты.}
\label{fig:vibemode}
\end{figure}

\textbf{Vibe Mode (WebGL2-визуализатор).} Полноэкранный визуализатор на fragment-шейдере GLSL~ES~3.00, реализующем SDF raymarching тоннеля с domain-repeated арками. Вход~--- 13~uniform-переменных, обновляемых через \code{requestAnimationFrame}: \code{uBass}/\code{uMid}/\code{uTreble} (FFT-полосы), \code{uBassEnv} (fast attack 0{,}70, slow release 0{,}12), \code{uLoudnessEnv}, \code{uOnset} (peak detector 600\,мс decay), \code{uArousal}/\code{uValence}, \code{uSegmentIntensity} (Intro=0{,}25, Chorus=1{,}0), \code{uSegmentFlash}, жанровые \code{uGenreHue}/\code{uGenreRibs}/\code{uGenreSharp}/\code{uGenreSparkle}/\code{uGenreTwist}. WebAudio: \code{fftSize=1024}, \code{smoothingTimeConstant=0.82}; чтение семантических данных с учётом \code{AudioContext.outputLatency} ($\approx$40--80\,мс).

\textbf{Унифицированный плеер.} Аудио-плеер и TimeAxis объединены в glass-карточку: фоновый waveform из \code{audio\_features.loudness\_rms} через \code{mix-blend-mode: screen}; монохромные segment-markers с \code{opacity=confidence} и uppercase-лейблами; glass-tooltip \code{T: 2:08 · ПРИПЕВ}; зелёный now-playing pulse.

\textbf{Обновлённые панели.} \emph{Эмоциональный профиль}~--- wide-span~4: A/V-trajectory с 4-квадрантной географией и хронологически окрашенным хвостом (индиго $\to$ магента $\to$ коралл $\to$ амбер), 4~hero-tile, pie сегментов, progress-бары. \emph{Динамика эмоций}~--- mood-tinted arousal + белая valence; rescale $[-1,1]\to[0,1]$ через $(x+1)\cdot 0{,}5$. \emph{Мел-спектрограмма}~--- aurora colormap + bloom; hover crosshair \code{T: 1:51 · F: 4.1\,кГц · I: 64\%}. \emph{Жанр}~--- $\star$~winner + \code{lighter} blend.

\textbf{Галерея треков.} Grid-галерея (Apple Music / Spotify style): 2/3/4/5 колонок, квадратные cover-tiles \code{aspect-ratio: 1/1}, fallback-градиент из \code{--track-mood $\times$ --track-tint}, duration chip, play-overlay при hover, Fraunces title.

\textbf{Извлечение метаданных.} Из ID3-тегов (npm \code{music-metadata}): artist, title, APIC~frame $\to$ \code{covers/\{id\}.jpg}. Из YouTube: \code{uploader} $\to$ artist, thumbnail $\to$ cover. Backfill при первом \code{GET /api/tracks}.

\textbf{Резюме главы 6.} Полнофункциональный веб-прототип с трёхзвенной архитектурой и синхронизированным шестипанельным дашбордом; итеративный редизайн (glass theme, Vibe~Mode, галерея, метаданные) превратил академический дашборд в продуктовый интерфейс.
